use case of pandas 

data clening 
data analysis 
data transformation
data visualization
data aggregation
file handling 
data filtering & selection 
time series analysis 



////////// PANDAS Series ///////////////    :- it is not importaant

% Pandas Series ek 1-dimensional labeled array hoti hai jisme data ke saath har value ka ek index (label) hota hai
1. "1-Dimensional" ka matlab kya hai?

% 1-Dimensional (1D) ka matlab hota hai ki data sirf ek line (ek column ya ek list) me store hota hai.
% [1,2,3,4,5] this is one dimantional array 
series vertically hota hai or array horigentaly hota hai 


2. "Array" ka matlab kya hai?

Array ka matlab hota hai same type ke multiple values ko ek saath store karna.


3. "Labeled" ka matlab kya hai?

Ye sabse important part hai.

Labeled ka matlab hai ki har value ke saath ek naam (label/index) hota hai.




How to create series 


import numpy as np
import pandas as pd




///////////////////////Creating a DataFrame///////////////////// 

DataFrame Kya Hai?

Definition:

Pandas DataFrame ek 2-dimensional labeled data structure hai, jisme data Rows aur Columns ke form me store hota hai.

Simple words me:

DataFrame ek Excel Sheet ya Database Table ki tarah hota hai.

Agar Series = Ek Column hai, to DataFrame = Multiple Columns hain.



///////////////////Selection And Indexing of columns///////////////////

Jab hum DataFrame bana lete hain, tab hume kisi particular column ya multiple columns ka data nikalna hota hai. Is process ko Selection of Columns kehte hain.

Q . Create a new column in list 
df["Designation"] = ["Doctor","Eng","Doctor","Eng"] 
df


///////////////////////Remove a column ////////////////////////////

# esse humlo column ko delete kr sakte hai 
df.drop("Designation", axis=1,inplace=True ) #implace ka matlab yah hai ki jab humlog kisi bhi chis ka drop krte hai toh usko inplace laga kr kisi bhi area ko remove kr sakte hai 

# Display DataFrame
print(df)



/////////////////////Selecting Rows///////////////////////////////////

Rows select karne ke liye Pandas me mainly 3 methods use hote hain:

1. loc[] → Label (Index Name) se row select karta hai.
2. iloc[] → Position (Row Number) se row select karta hai.
3. Condition (Filtering) → Kisi condition ke basis par rows select karta hai.


\\\\\\\\\\\\\\\\\\\\ Conditional Selection////////////////////////

Conditional Selection ka matlab hai kisi condition ke basis par rows ya data ko select karna.

Jaise:

Marks 80 se zyada wale students.
Age 18 se kam wale students.
Salary 50000 se zyada wale employees.

Ye sab Conditional Selection ke examples hain.




///////////////////////////////////     Missing Data      /////////////////////////////

Finding missing data 

Missing Data ka matlab hota hai ki dataset me kuch values available nahi hoti ya blank (NaN/Null) hoti hain. Ye data collection ya data entry ke time kisi wajah se missing ho sakti hain.

Types of Missing Data
1. Missing Completely at Random (MCAR)
Missing value ka kisi bhi dusre data se koi relation nahi hota.
Ye bilkul random hota hai.



job hum do data ko murge krna chate hai toh hum 

pd.merge(data1 name, data2 name , on = 'emoployee_id', how = "inner')


//////////////// Thresh old ka matlab hota hai ki humlog kitne missing value ko allow karenge kisi bhi row me ya column me. Agar humlog 2 ka threshold set karte hai toh iska matlab hai ki agar kisi bhi row me 2 se zyada missing value hai toh usko drop kar do otherwise usko rakh do.////////////////



